\chapter{فضای برداری}
\label{ch:vector-spaces}

تا اینجا با بردارهای $\R^n$ کار کردیم: آن‌ها را جمع می‌زدیم و در عدد ضرب می‌کردیم. اما ماتریس‌ها هم همین رفتار را دارند؛ چندجمله‌ای‌ها و توابع هم! دو چندجمله‌ای را می‌شود جمع زد و یک چندجمله‌ای را می‌شود در عدد ضرب کرد و همان قواعد آشنا برقرار است. ایده‌ی \textbf{فضای برداری} این است که این قواعد مشترک را جدا کنیم و یک بار برای همیشه درباره‌ی هر چیزی که از آن‌ها پیروی می‌کند حرف بزنیم. فضای برداری \emph{سنگ‌بنای} جبر خطی است.

\section{عمل دوتایی}

\begin{definition}[عمل دوتایی]
هر تابع $*:A\times A\to A$ یک \term{عمل دوتایی}{Binary Operation} روی مجموعه‌ی $A$ است؛ یعنی قاعده‌ای که به هر دو عضو $a,b\in A$ یک عضو $a*b$ را \textbf{از خود $A$} نسبت می‌دهد.
\end{definition}

\begin{definition}[بسته بودن]
مجموعه‌ی $A$ تحت یک عمل (مثل جمع، تفریق، ضرب و …) \term{بسته}{Closed} است اگر حاصل انجام آن عمل روی اعضای $A$ همیشه دوباره در $A$ باشد:
\[
a\in A,\ b\in A\ \Longrightarrow\ a*b\in A .
\]
به این ویژگی \term{قانون بستاری}{Closure Law} می‌گوییم.
\end{definition}

\begin{example}[بسته هست یا نه؟]
\begin{itemize}
  \item $\N$ تحت جمع و ضرب بسته است، اما تحت تفریق نه: $2-3=-1\notin\N$.
  \item $\Z$ تحت جمع، تفریق و ضرب بسته است، اما تحت تقسیم نه: $1\div2\notin\Z$.
  \item مجموعه‌ی اعداد فرد تحت جمع بسته نیست: $1+1=2$.
  \item مجموعه‌ی ماتریس‌های $n\times n$ تحت ضرب ماتریسی بسته است؛ اما مجموعه‌ی \emph{همه‌ی} ماتریس‌ها نه، چون ضرب همه‌ی جفت‌ها تعریف نشده است.
\end{itemize}
\end{example}

\begin{note}
وقتی می‌گوییم «$*$ یک عمل دوتایی روی $A$ است»، بسته بودن $A$ تحت $*$ را هم به‌طور ضمنی گفته‌ایم؛ چون خروجی تابع باید در $A$ باشد.
\end{note}

\section{تعریف فضای برداری}

یک فضای برداری از سه چیز ساخته می‌شود: مجموعه‌ای ناتهی $V$ (که اعضایش را «بردار» می‌نامیم)، یک میدان $\F$ (که اعضایش را «اسکالر» می‌نامیم؛ بیشتر وقت‌ها $\R$ یا $\C$) و دو عمل:
\begin{enumerate}
  \item \textbf{جمع برداری} $+:V\times V\to V$؛ یک عمل دوتایی روی $V$.
  \item \textbf{ضرب اسکالر} $\cdot:\F\times V\to V$؛ که یک اسکالر و یک بردار می‌گیرد و یک \emph{بردار} برمی‌گرداند.
\end{enumerate}

\begin{note}
در این درس، اگر چیزی نگوییم، میدان به‌طور پیش‌فرض $\R$ (اعداد حقیقی) است.
\end{note}

\begin{definition}[فضای برداری]
یک \term{فضای برداری}{Vector Space} روی میدان $\F$، مجموعه‌ی ناتهی $V$ همراه با دو عمل است که آن را با $(V,\F,+,\cdot)$ نشان می‌دهیم:
\begin{enumerate}[label=(\roman*)]
  \item \textbf{جمع برداری}: به هر $\vx,\vy\in V$ عضو $\vx+\vy\in V$ را نسبت می‌دهد.
  \item \textbf{ضرب اسکالر}: به هر $\alpha\in\F$ و $\vx\in V$ عضو $\alpha\cdot\vx\in V$ را نسبت می‌دهد. معمولاً نقطه را نمی‌نویسیم: $\alpha\vx$.
\end{enumerate}
و این هشت اصل برای همه‌ی $\vx,\vy,\vz\in V$ و $\alpha,\beta\in\F$ برقرار باشد:

\smallskip
\textbf{اصول جمع:}
\begin{enumerate}
  \item \textbf{جابه‌جایی:} $\vx+\vy=\vy+\vx$.
  \item \textbf{شرکت‌پذیری:} $(\vx+\vy)+\vz=\vx+(\vy+\vz)$.
  \item \textbf{عضو خنثی (بردار صفر):} برداری مثل $\zero\in V$ وجود دارد که برای هر $\vx$: $\vx+\zero=\vx$.
  \item \textbf{قرینه:} برای هر $\vx$ برداری مثل $-\vx\in V$ وجود دارد که $\vx+(-\vx)=\zero$.
\end{enumerate}
\textbf{اصول ضرب اسکالر (اثر میدان روی فضا):}
\begin{enumerate}[resume]
  \item \textbf{شرکت‌پذیری:} $\alpha(\beta\vx)=(\alpha\beta)\vx$.
  \item \textbf{توزیع‌پذیری روی جمع اسکالرها:} $(\alpha+\beta)\vx=\alpha\vx+\beta\vx$.
  \item \textbf{توزیع‌پذیری روی جمع بردارها:} $\alpha(\vx+\vy)=\alpha\vx+\alpha\vy$.
  \item \textbf{همانی اسکالر:} $1\vx=\vx$.
\end{enumerate}
\end{definition}

\begin{intuition}
به این اصول مثل «قوانین بازی» نگاه کنید. لازم نیست بردارها پیکان باشند؛ هر چیزی که بتوان آن را جمع زد و در عدد ضرب کرد و این هشت قانون را رعایت کند، بردار است. از این به بعد هر قضیه‌ای که فقط با این هشت اصل ثابت کنیم، هم‌زمان برای $\R^n$، ماتریس‌ها، چندجمله‌ای‌ها و توابع درست است.
\end{intuition}

\begin{warning}
در یک فضای برداری کلی، \textbf{ضرب دو بردار در هم} تعریف نشده است و طول و زاویه هم (هنوز) معنی ندارد. فقط دو کار مجاز است: جمع بردارها و ضرب بردار در اسکالر. مفهوم طول و زاویه را در فصل ۸ با ضرب داخلی اضافه می‌کنیم.
\end{warning}

\section{نتایج اولیه‌ی اصول}

چند ویژگی که «بدیهی» به نظر می‌رسند در واقع جزو اصول نیستند و باید از آن‌ها ثابت شوند. این اثبات‌ها تمرین خوبی برای کار با اصول هستند.

\begin{theorem}[ویژگی‌های پایه]\label{thm:vs-basic}
در هر فضای برداری $V$ روی $\F$، برای هر $\vx,\vy,\vz\in V$ و $\alpha\in\F$:
\begin{enumerate}
  \item \textbf{قانون حذف:} اگر $\vx+\vy=\vx+\vz$ آن‌گاه $\vy=\vz$.
  \item بردار صفر یکتاست.
  \item قرینه‌ی هر بردار یکتاست.
  \item $0\vx=\zero$ \ (اسکالر صفر ضرب در هر بردار، بردار صفر است).
  \item $\alpha\zero=\zero$.
  \item $(-1)\vx=-\vx$.
  \item اگر $\alpha\vx=\zero$ آن‌گاه $\alpha=0$ یا $\vx=\zero$.
\end{enumerate}
\end{theorem}
\begin{proof}
در هر قدم، اصلی را که استفاده می‌کنیم در پرانتز می‌نویسیم.

\textbf{۱.} به دو طرف $\vx+\vy=\vx+\vz$ بردار $-\vx$ را از چپ اضافه می‌کنیم:
\[
\vy\overset{(3)}{=}\zero+\vy\overset{(4)}{=}\big((-\vx)+\vx\big)+\vy\overset{(2)}{=}(-\vx)+(\vx+\vy)=(-\vx)+(\vx+\vz)\overset{(2)}{=}\big((-\vx)+\vx\big)+\vz\overset{(4)}{=}\zero+\vz\overset{(3)}{=}\vz .
\]
(از جابه‌جایی هم استفاده کردیم تا $\zero+\vy=\vy+\zero$ و $(-\vx)+\vx=\vx+(-\vx)$ باشد.)

\textbf{۲.} اگر $\zero$ و $\zero'$ هر دو خاصیت عضو خنثی را داشته باشند:
$\zero'=\zero'+\zero=\zero+\zero'=\zero$.
(تساوی اول چون $\zero$ خنثی است، دومی جابه‌جایی، سومی چون $\zero'$ خنثی است.)

\textbf{۳.} اگر $\vy$ و $\vz$ هر دو قرینه‌ی $\vx$ باشند، $\vx+\vy=\zero=\vx+\vz$ و طبق قانون حذف $\vy=\vz$.

\textbf{۴.} با توزیع‌پذیری روی جمع اسکالرها:
$0\vx=(0+0)\vx=0\vx+0\vx$.
پس $0\vx+\zero=0\vx+0\vx$ و طبق قانون حذف $0\vx=\zero$.

\textbf{۵.} با توزیع‌پذیری روی جمع بردارها:
$\alpha\zero=\alpha(\zero+\zero)=\alpha\zero+\alpha\zero$
و مثل قبل $\alpha\zero=\zero$.

\textbf{۶.} داریم
$\vx+(-1)\vx=1\vx+(-1)\vx=(1+(-1))\vx=0\vx=\zero$
(اصل ۸، اصل ۶ و بند ۴). پس $(-1)\vx$ یک قرینه‌ی $\vx$ است و چون قرینه یکتاست، $(-1)\vx=-\vx$.

\textbf{۷.} فرض کنید $\alpha\vx=\zero$ و $\alpha\neq0$. چون $\F$ میدان است، $\alpha^{-1}$ وجود دارد و:
\[
\vx=1\vx=(\alpha^{-1}\alpha)\vx=\alpha^{-1}(\alpha\vx)=\alpha^{-1}\zero=\zero .
\]
\end{proof}

\begin{note}
بند ۷ جایی است که «میدان» بودن $\F$ (وجود وارون ضربی) واقعاً به کار می‌آید. تفریق را هم این‌طور تعریف می‌کنیم: $\vx-\vy:=\vx+(-\vy)$.
\end{note}

\section{نمونه‌هایی از فضای برداری}

\subsection{فضای $\R^n$ (و $\F^n$)}
$\R^n$ با جمع و ضرب اسکالرِ مؤلفه‌به‌مؤلفه یک فضای برداری روی $\R$ است. هر هشت اصل مستقیماً از ویژگی‌های اعداد حقیقی می‌آیند، چون همه‌چیز مؤلفه‌به‌مؤلفه انجام می‌شود. مثلاً جابه‌جایی:
\[
\vx+\vy=\begin{bmatrix}x_1+y_1\\ \vdots\\ x_n+y_n\end{bmatrix}=\begin{bmatrix}y_1+x_1\\ \vdots\\ y_n+x_n\end{bmatrix}=\vy+\vx .
\]
بردار صفر $\zero=(0,\dots,0)$ و قرینه‌ی $\vx$ برابر $(-x_1,\dots,-x_n)$ است. به همین ترتیب برای هر میدان $\F$، مجموعه‌ی $n$تایی‌های $\F^n$ روی $\F$ فضای برداری است (مثلاً $\C^n$ روی $\C$).

\subsection{فضای ماتریس‌ها $M_{m,n}(\R)$}
مجموعه‌ی همه‌ی ماتریس‌های $m\times n$ با درایه‌های حقیقی، با جمع ماتریسی و ضرب اسکالر معمولی، یک فضای برداری روی $\R$ است. دلیلش دقیقاً مثل $\R^n$ است: هر ماتریس $m\times n$ در واقع $mn$ عدد است که درایه‌به‌درایه جمع و ضرب می‌شوند. بردار صفر، ماتریس صفر است.

\subsection{فضای توابع}
فرض کنید $X$ یک مجموعه‌ی ناتهی و $\F$ یک میدان باشد. مجموعه‌ی همه‌ی توابع از $X$ به $\F$ را با $\F^X$ نشان می‌دهیم:
\[
\F^X=\{f\ :\ f:X\to\F\}.
\]
جمع و ضرب اسکالر را \textbf{نقطه‌به‌نقطه} تعریف می‌کنیم:
\[
(f+g)(x):=f(x)+g(x),\qquad (af)(x):=a\,f(x)\qquad (x\in X).
\]
یعنی تابع $f+g$ در هر نقطه مقدارش مجموع مقادیر $f$ و $g$ در آن نقطه است.

\begin{theorem}
$\F^X$ با جمع و ضرب اسکالر نقطه‌به‌نقطه، یک فضای برداری روی $\F$ است.
\end{theorem}
\begin{proof}
یادآوری: دو تابع برابرند اگر در همه‌ی نقاط مقدار برابر داشته باشند. پس برای اثبات هر اصل، کافی است تساوی را در یک نقطه‌ی دلخواه $x\in X$ بررسی کنیم؛ و در آن نقطه همه‌چیز به ویژگی‌های میدان $\F$ برمی‌گردد. مثلاً:
\begin{itemize}
  \item \textbf{جابه‌جایی:} $(f+g)(x)=f(x)+g(x)=g(x)+f(x)=(g+f)(x)$.
  \item \textbf{بردار صفر:} \term{تابع صفر}{Zero Function} $\mathbf{0}(x):=0$ برای همه‌ی $x$. داریم $(f+\mathbf{0})(x)=f(x)+0=f(x)$.
  \item \textbf{قرینه:} $(-f)(x):=-f(x)$. آن‌گاه $(f+(-f))(x)=f(x)-f(x)=0=\mathbf{0}(x)$.
  \item \textbf{توزیع‌پذیری:} $\big(a(f+g)\big)(x)=a\big(f(x)+g(x)\big)=af(x)+ag(x)=(af+ag)(x)$.
\end{itemize}
بقیه‌ی اصول هم دقیقاً به همین شکل ثابت می‌شوند.
\end{proof}

\begin{example}[چند فضای تابعی مهم]
\begin{itemize}
  \item $\R^{\R}$: همه‌ی توابع حقیقی روی $\R$.
  \item $P(\R)$: مجموعه‌ی همه‌ی چندجمله‌ای‌ها با ضرایب حقیقی.
  \item $C[0,1]$: مجموعه‌ی همه‌ی توابع حقیقیِ \textbf{پیوسته} روی $[0,1]$.
  \item مجموعه‌ی همه‌ی توابع حقیقیِ \textbf{مشتق‌پذیر} روی $[0,1]$.
\end{itemize}
همه‌ی این‌ها با عمل‌های نقطه‌به‌نقطه فضای برداری‌اند. برای سه مورد آخر کافی است بدانیم «مجموع دو چندجمله‌ای/پیوسته/مشتق‌پذیر، باز از همان نوع است و ضرب آن در عدد هم» (بستاری)؛ بقیه‌ی اصول از $\R^{\R}$ یا $\R^{[0,1]}$ به ارث می‌رسند. (این ایده را در فصل بعد با مفهوم \emph{زیرفضا} دقیق می‌کنیم.) حتی مجموعه‌ی همه‌ی \textbf{دنباله‌های} حقیقی $(a_1,a_2,a_3,\dots)$ هم فضای برداری است؛ چون یک دنباله چیزی جز تابعی از $\N$ به $\R$ نیست.
\end{example}

\subsection{فضای چندجمله‌ای‌ها $P_n(\R)$}
$P_n(\R)$ مجموعه‌ی همه‌ی چندجمله‌ای‌ها با ضرایب حقیقی و \textbf{درجه‌ی حداکثر $n$} است:
\[
p(x)=a_0+a_1x+a_2x^2+\dots+a_nx^n,\qquad a_i\in\R .
\]
\begin{itemize}
  \item \textbf{جمع:} ضرایب هم‌درجه جمع می‌شوند:
  $(p+q)(x)=(a_0+b_0)+(a_1+b_1)x+\dots+(a_n+b_n)x^n$.
  \item \textbf{ضرب اسکالر:} $(cp)(x)=ca_0+ca_1x+\dots+ca_nx^n$.
\end{itemize}
هر دو حاصل باز چندجمله‌ای با درجه‌ی حداکثر $n$ هستند و هشت اصل دیگر هم برقرارند (درست مثل $\R^{n+1}$؛ در واقع چندجمله‌ای $p$ را می‌شود با $n+1$ ضریبش، $(a_0,\dots,a_n)$، یکی گرفت). بردار صفر، چندجمله‌ای صفر است (همه‌ی ضرایب صفر).

\begin{warning}
مجموعه‌ی چندجمله‌ای‌های با درجه‌ی \textbf{دقیقاً} $n$ فضای برداری \emph{نیست}: اولاً چندجمله‌ای صفر در آن نیست و ثانیاً تحت جمع بسته نیست؛ مثلاً $(x^2+x)+(-x^2)=x$ که درجه‌اش ۱ است.
\end{warning}

\subsection{چند مثال و مثال نقض}
\begin{example}[کدام‌ها فضای برداری‌اند؟]
با جمع و ضرب معمولی، کدام یک فضای برداری است؟
\begin{multicols}{2}
\begin{enumerate}
  \item $\R^n$ روی $\R$
  \item $\C$ روی $\R$
  \item $\R$ روی $\C$
  \item $\Z$ روی $\R$
  \item همه‌ی چندجمله‌ای‌ها با ضرایب حقیقی، روی $\R$
  \item چندجمله‌ای‌های درجه‌ی حداکثر $n$، روی $\R$
  \item ماتریس‌ها $M_{m,n}(\R)$ روی $\R$
  \item توابع $f:X\to\R$ روی $\R$
\end{enumerate}
\end{multicols}
\end{example}
\begin{solution}
\begin{enumerate}
  \item \textbf{بله} (بخش قبل).
  \item \textbf{بله.} جمع دو عدد مختلط، مختلط است و ضرب عدد حقیقی در عدد مختلط هم مختلط است؛ اصول از ویژگی‌های اعداد مختلط می‌آیند. (در واقع $a+bi$ مثل بردار $(a,b)\in\R^2$ رفتار می‌کند.)
  \item \textbf{خیر.} ضرب اسکالر بسته نیست: اسکالر $i\in\C$ ضرب در «بردار» $1\in\R$ می‌شود $i\notin\R$.
  \item \textbf{خیر.} ضرب اسکالر بسته نیست: $\frac12\cdot1=\frac12\notin\Z$.
  \item \textbf{بله} ($P(\R)$).
  \item \textbf{بله} ($P_n(\R)$).
  \item \textbf{بله.}
  \item \textbf{بله} ($\R^X$).
\end{enumerate}
\end{solution}

\begin{example}[عمل‌های عجیب]
فرض کنید $V=\R$ (همه‌ی اعداد حقیقی) با این عمل‌ها باشد:
\[
u\oplus v:=u-v\quad(\text{\fa{تفریق معمولی}}),\qquad \alpha\odot u:=\alpha u\quad(\text{\fa{ضرب معمولی}}).
\]
آیا $V$ فضای برداری است؟ اگر نه، کدام اصول برقرار نیستند؟
\end{example}
\begin{solution}
بستاری برقرار است (تفریق و ضرب دو عدد حقیقی، حقیقی است). حالا اصول را یکی‌یکی بررسی می‌کنیم:
\begin{enumerate}
  \item \textbf{جابه‌جایی: برقرار نیست.} $u\oplus v=u-v$ اما $v\oplus u=v-u$؛ مثلاً $1\oplus0=1\neq-1=0\oplus1$.
  \item \textbf{شرکت‌پذیری: برقرار نیست.} $(u\oplus v)\oplus w=u-v-w$ اما $u\oplus(v\oplus w)=u-v+w$؛ مثلاً با $u=v=0,\ w=1$: $-1\neq1$.
  \item \textbf{بردار صفر: به همان شکلی که در تعریف نوشتیم برقرار است.} با عدد $0$ داریم $u\oplus0=u-0=u$. البته $0$ از چپ خنثی نیست ($0\oplus u=-u$)؛ در فضاهای برداری واقعی، جابه‌جایی این مشکل را برطرف می‌کند اما این‌جا جابه‌جایی نداریم.
  \item \textbf{قرینه: برقرار است.} $u\oplus u=u-u=0$؛ پس هر عضو قرینه‌ی خودش است.
  \item \textbf{شرکت‌پذیری ضرب: برقرار است.} $\alpha\odot(\beta\odot u)=\alpha\beta u=(\alpha\beta)\odot u$.
  \item \textbf{توزیع روی جمع اسکالرها: برقرار نیست.} $(\alpha+\beta)\odot u=(\alpha+\beta)u$ اما $(\alpha\odot u)\oplus(\beta\odot u)=\alpha u-\beta u$؛ مثلاً $\alpha=\beta=1,\ u=1$: $2\neq0$.
  \item \textbf{توزیع روی جمع بردارها: برقرار است.} $\alpha\odot(u\oplus v)=\alpha(u-v)=\alpha u-\alpha v=(\alpha\odot u)\oplus(\alpha\odot v)$.
  \item \textbf{همانی: برقرار است.} $1\odot u=u$.
\end{enumerate}
پس $V$ فضای برداری \textbf{نیست}؛ اصول جابه‌جایی، شرکت‌پذیری جمع و توزیع‌پذیری روی جمع اسکالرها نقض می‌شوند.
\end{solution}

\begin{example}[یک فضای برداری غیرمنتظره]
فرض کنید $V=\R^{+}$ (اعداد حقیقی \emph{مثبت}) باشد با عمل‌های
\[
u\oplus v:=uv\quad(\text{\fa{ضرب معمولی}}),\qquad \alpha\odot u:=u^{\alpha}\quad(\text{\fa{توان}}).
\]
نشان دهید $V$ یک فضای برداری روی $\R$ است.
\end{example}
\begin{solution}
بستاری: حاصل‌ضرب دو عدد مثبت و توان یک عدد مثبت، مثبت است. اصول:
\begin{itemize}
  \item جابه‌جایی و شرکت‌پذیری: $uv=vu$ و $(uv)w=u(vw)$.
  \item \textbf{بردار صفر عدد ۱ است!} چون $u\oplus1=u\cdot1=u$.
  \item قرینه‌ی $u$ عدد $\frac1u$ است: $u\oplus\frac1u=1$ (که همان «صفرِ» این فضاست).
  \item $\alpha\odot(\beta\odot u)=(u^\beta)^\alpha=u^{\alpha\beta}=(\alpha\beta)\odot u$.
  \item $(\alpha+\beta)\odot u=u^{\alpha+\beta}=u^\alpha u^\beta=(\alpha\odot u)\oplus(\beta\odot u)$.
  \item $\alpha\odot(u\oplus v)=(uv)^\alpha=u^\alpha v^\alpha=(\alpha\odot u)\oplus(\alpha\odot v)$.
  \item $1\odot u=u^1=u$.
\end{itemize}
پس همه‌ی اصول برقرارند. نکته‌ی آموزنده: «بردار صفر» لزوماً عدد صفر نیست؛ هر عضوی است که نقش خنثیِ \emph{جمعِ آن فضا} را بازی کند.
\end{solution}

\begin{note}[کوچک‌ترین فضای برداری]
مجموعه‌ی تک‌عضوی $\{\zero\}$ با $\zero+\zero=\zero$ و $\alpha\zero=\zero$ یک فضای برداری است که به آن \term{فضای صفر}{Zero Space} می‌گوییم.
\end{note}

\section{جمع‌بندی و تمرین}

\begin{summary}
\begin{itemize}
  \item فضای برداری $(V,\F,+,\cdot)$: مجموعه‌ای که بستاری تحت جمع و ضرب اسکالر دارد و هشت اصل (چهار اصل جمع و چهار اصل ضرب اسکالر) در آن برقرار است.
  \item از اصول نتیجه می‌شود: صفر و قرینه یکتا هستند، $0\vx=\zero$، $\alpha\zero=\zero$، $(-1)\vx=-\vx$ و $\alpha\vx=\zero\Rightarrow\alpha=0\ \text{\fa{یا}}\ \vx=\zero$.
  \item مثال‌های اصلی: $\R^n$ و $\F^n$، ماتریس‌ها $M_{m,n}$، توابع $\F^X$ (با عمل‌های نقطه‌به‌نقطه)، چندجمله‌ای‌ها $P(\R)$ و $P_n(\R)$، توابع پیوسته و مشتق‌پذیر.
  \item برای رد کردن فضای برداری بودن، کافی است یک اصل (یا بستاری) با یک مثال نقض شکسته شود.
\end{itemize}
\end{summary}

\begin{exercise}
آیا $\{(x,y)\in\R^2: x\ge0\}$ با عمل‌های معمولی فضای برداری است؟
\end{exercise}

\begin{exercise}
روی $\R^2$ جمع معمولی را در نظر بگیرید اما ضرب اسکالر را این‌طور تعریف کنید: $c\odot(x,y)=(cx,0)$. کدام اصل برقرار نیست؟
\end{exercise}

\begin{exercise}
فقط با اصول فضای برداری ثابت کنید $-(-\vx)=\vx$.
\end{exercise}

\begin{exercise}
ثابت کنید اگر $\alpha\vx=\beta\vx$ و $\vx\neq\zero$، آن‌گاه $\alpha=\beta$.
\end{exercise}

\begin{exercise}
آیا مجموعه‌ی ماتریس‌های $2\times2$ که مجموع درایه‌هایشان ۱ است، با عمل‌های معمولی فضای برداری است؟
\end{exercise}

\subsection*{پاسخ‌های کوتاه}
\begin{enumerate}
  \item خیر؛ ضرب اسکالر بسته نیست: $(-1)(1,0)=(-1,0)$ در مجموعه نیست.
  \item فقط اصل همانی: $1\odot(x,y)=(x,0)\neq(x,y)$ وقتی $y\neq0$. (بقیه‌ی اصول برقرارند؛ پس اصل همانی واقعاً لازم است و از بقیه نتیجه نمی‌شود.)
  \item $-\vx$ قرینه‌ی $\vx$ است: $(-\vx)+\vx=\zero$. پس $\vx$ یک قرینه‌ی $-\vx$ است و چون قرینه یکتاست، $-(-\vx)=\vx$.
  \item $(\alpha-\beta)\vx=\alpha\vx-\beta\vx=\zero$ و چون $\vx\neq\zero$، طبق بند ۷ قضیه‌ی \ref{thm:vs-basic}، $\alpha-\beta=0$.
  \item خیر؛ ماتریس صفر در آن نیست (و تحت جمع هم بسته نیست: مجموع درایه‌های حاصل‌جمع ۲ می‌شود).
\end{enumerate}
